\documentclass[11pt]{article}
\usepackage[margin=1in]{geometry}
\usepackage{amsmath,amssymb,amsthm,booktabs}
\usepackage[hidelinks]{hyperref}
\usepackage{xurl}
\setlength{\emergencystretch}{3em}
\newtheorem{theorem}{Theorem}
\newtheorem{lemma}[theorem]{Lemma}
\newtheorem{corollary}[theorem]{Corollary}
\newcommand{\Z}{\mathbb Z}
\newcommand{\F}{\mathbb F}
\newcommand{\A}{\mathcal A}
\title{Uniform Fox-ideal model barriers\protect\\for a prime four-bridge family}
\author{}
\date{18 November 2026}
\begin{document}\maketitle
\begin{abstract}
For a fixed four-meridian presentation of an explicit prime twist
family, the abelianized Fox ideal is the whole integer Laurent ring.
An explicit identity also makes its image the whole group algebra
$\F_2[A_8]$ under a fixed meridional quotient. Both identities use only
the first relator, which is unchanged by the entire twist operation.
Consequently scalar coefficient models fail over every commutative
ring, and all characteristic-two models through this particular
quotient fail in every dimension. An integer determinant argument
also excludes every unital characteristic-zero model through this quotient,
leaving only finitely many possible odd characteristics. The exclusion
extends to nilpotent lifts of characteristic two and to local models
with this residue quotient. These are restrictions on the Fox-ideal rank
criterion; the ordinary knot-group rank remains undetermined.
\end{abstract}

Put
\begin{equation}\label{eq:family}
 U=\sigma_1\sigma_2\sigma_3^{-2}\sigma_2\sigma_3^{-2}
 \sigma_2\sigma_1^2\sigma_2^3\sigma_3^{-1},\qquad
 K_n=\widehat{U\sigma_2^{3+6n}}\quad(n\ge0).
\end{equation}
These exact knots have already been proved prime with
$\operatorname{mr}(K_n)=b(K_n)=4$ and complex augmentation rank three
\cite{family}. That result motivates the rank search; the algebraic
identities below do not require its geometric or augmentation imports.

For a deficiency-one presentation
$G=\langle x_0,\ldots,x_{g-1}\mid R_0,\ldots,R_{g-2}\rangle$,
let $I_X\subset\Z[G]$ be the two-sided ideal generated by all images
of the left Fox derivatives $\partial R_j/\partial x_i$.
The rank criterion stated by Moriah \cite[pp.~197--198, Theorem~3.5]{moriah}
says that a nontrivial representation
$\Z[G]/I_X\longrightarrow M_m(R)$, with $R$ a commutative ring with
unit, forces $\operatorname{rank}(G)=g$.
We examine particular possible coefficient models for this criterion.
The original Lustig--Moriah lemma underlying the cited survey theorem
is an imported result, not an independently checked proof here.
No converse to the rank criterion is used.

\section{A literal presentation and a fixed relator}
Let $F=F(x_0,x_1,x_2,x_3)$. Propagate ordered free words along a
braid, starting with $(x_0,x_1,x_2,x_3)$, by the rules
\begin{equation}\label{eq:action}
 (a,b)\stackrel{\sigma_i}{\longmapsto}(b,b^{-1}ab),\qquad
 (a,b)\stackrel{\sigma_i^{-1}}{\longmapsto}(aba^{-1},a)
\end{equation}
on slots $i-1,i$. Write $F_i^{(n)}$ for the four outputs of
$U\sigma_2^{3+6n}$ and $R_i^{(n)}=F_i^{(n)}x_i^{-1}$.

\begin{lemma}\label{lem:presentation}
The ordinary knot group has presentation
\[
 G_n=\pi_1(S^3\setminus K_n)
   =\langle x_0,x_1,x_2,x_3\mid R_0^{(n)},R_1^{(n)},R_2^{(n)}\rangle.
\]
The relator $R_0^{(n)}$ is one identical free word $R_0$ for all $n$.
\end{lemma}
\begin{proof}
The rules \eqref{eq:action} are the oriented Wirtinger relations
solved for the outgoing underarc. Eliminating internal arc generators
and identifying matching closure endpoints gives all four relations
$F_i^{(n)}=x_i$. Each local rule preserves the ordered product $ab$,
so $F_0^{(n)}F_1^{(n)}F_2^{(n)}F_3^{(n)}=x_0x_1x_2x_3$ in $F$.
The first three closure relations therefore imply the fourth.
The braid permutation at $n=0$ is a four-cycle, and the extra $6n$
transpositions preserve it. Thus every closure is a knot, and these
four generators are conjugate oriented meridians.

Each terminal $\sigma_2$ changes slots $1,2$ and leaves slot $0$
unchanged as a free word. Induction on its exponent proves
$F_0^{(n)}=F_0^{(0)}$ and hence the asserted equality of relators.
For the frozen signed diagram of $K_0$, the ordered top arcs are
$(1,0,8,13)$. The exact checker follows all seventeen signed crossings
with these labels and recovers \eqref{eq:action} and closure.
It also verifies the ordered-product identity by free reduction.
This is a literal diagram binding, rather than an identification of
unrelated simplified generators. The first relator has length $156$.
\end{proof}

Write $D_{ji}^{(n)}=\partial R_j^{(n)}/\partial x_i$ in $\Z[G_n]$,
and put $I_n=(D_{ji}^{(n)}:0\le j\le2,\ 0\le i\le3)$.
Our left derivative convention is
\[
 \partial_i(vw)=\partial_i(v)+v\partial_i(w),\qquad
 \partial_i(x_i^{-1})=-x_i^{-1}.
\]
Every coefficient computation below comes directly from this rule.

\section{An integer Laurent unit identity}
Let $H_{ji}^{(n)}(t)$ be the image of $D_{ji}^{(n)}$ under the
abelianization $x_i\mapsto t$ into $\Lambda=\Z[t,t^{-1}]$.

\begin{theorem}\label{thm:scalar}
For every $n\ge0$, the image of $I_n$ in $\Lambda$ is the unit ideal.
There is no nontrivial scalar ring representation of $\Z[G_n]/I_n$
into any commutative ring.
\end{theorem}
\begin{proof}
By Lemma~\ref{lem:presentation}, the first row $H_{0i}^{(n)}=H_{0i}$
is independent of $n$. Exact differentiation of $R_0$ gives
\begin{align*}
 f(t):=t^4H_{01}(t)&=t^7-4t^6+8t^5-12t^4+15t^3-11t^2+5t-1,\\
 g(t):=t^3H_{02}(t)&=t^7-4t^6+8t^5-11t^4+12t^3-9t^2+4t-1.
\end{align*}
Define
\begin{align*}
 A(t)&=-3t^6+5t^5+t^4-10t^3+15t^2-23t+16,\\
 B(t)&=3t^6-5t^5-t^4+7t^3-13t^2+35t-17.
\end{align*}
Integer coefficient multiplication gives
\begin{equation}\label{eq:scalarunit}
                  A(t)f(t)+B(t)g(t)=1.
\end{equation}
Thus $t^4AH_{01}+t^3BH_{02}=1$ in $\Lambda$.
This is an integer identity, with no rational denominator or inference
from a greatest common divisor over $\mathbb Q$.

In a unital scalar representation into a commutative ring, conjugate
meridians have the same unit image $u$. The representation factors
through $\Lambda$ by $t\mapsto u$. Killing all Fox entries would send
\eqref{eq:scalarunit} to $1=0$. For a nonzero ring homomorphism not
assumed unital, put $e=\rho(1)$. Its image lies in the nonzero
commutative ring $eR$ with identity $e$, and the same argument forces
$e=0$. This also excludes that convention.
\end{proof}

\section{A unit in the full characteristic-two group algebra}
Use permutations of $\{0,\ldots,7\}$, composed by
$(pq)(i)=p(q(i))$. The fixed top meridian images are
\begin{equation}\label{eq:top}
 p_0=(01)(23),\quad p_1=(14)(23),\quad
 p_2=(05)(16),\quad p_3=(12)(67).
\end{equation}
They give the same quotient $q_n:G_n\twoheadrightarrow A_8$ for all $n$.
Indeed, the exact seventeen-crossing propagation returns these images
at $n=0$. The two middle images after that return are $p_1,p_2$.
Direct permutation substitution in \eqref{eq:action} shows that six
further $\sigma_2$ steps return this pair. Induction gives closure
for all $n$, while exact enumeration of the subgroup generated by
\eqref{eq:top} gives all $20160$ even permutations.
Let $d_i\in\F_2[A_8]$ be the image of $D_{0i}^{(n)}$.
Both the free word and the four top images are fixed, so $d_i$ does
not depend on $n$.

\begin{theorem}\label{thm:a8}
For every $n\ge0$, the image of $I_n$ under
$\Z[G_n]\longrightarrow\F_2[A_8]$ is the whole ring.
Consequently no nonzero characteristic-two coefficient model
factoring through this particular quotient can kill all the Fox
entries, regardless of dimension or semisimplicity.
\end{theorem}
\begin{proof}
The accompanying exact certificate specifies elements $X_0,X_1$ of
$\F_2[A_8]$, with respective supports of sizes $9792$ and $369$,
and verifies the identity
\begin{equation}\label{eq:a8unit}
                         X_0d_0+X_1d_1=1.
\end{equation}
The complete finite data are in
\texttt{uniform\_fox\_a8\_char2\_certificate.json}.
Each support member is given as an explicit eight-entry permutation
tuple, and every nonzero coefficient is one. The checker independently
reconstructs the derivatives by their prefix products, validates every
permutation and sums the products modulo two; the only surviving
permutation is the identity. Thus verification needs neither a
regular-matrix rank computation nor a classification of simple modules.
Details of the exact replay are specified below.

Equation~\eqref{eq:a8unit} puts one in the image ideal. Any unital ring
homomorphism from $\F_2[A_8]$ killing the derivatives would then kill
one. This includes every action on a nonzero vector space, every
matrix dimension and every commutative characteristic-two coefficient
ring, with no complete-reducibility assumption.
Lemma~\ref{lem:presentation} makes this same finite identity valid
for every $n$; no sample-to-family inference is made.
\end{proof}

\begin{corollary}\label{cor:characteristic}
There exists a nonzero odd integer $d$ such that $d\cdot1$ belongs to
the image of $I_n$ in $\Z[A_8]$, for every $n$.
A unital model through $q_n$ killing $I_n$ can therefore have only
positive odd coefficient-ring characteristic dividing $|d|$.
In particular every unital characteristic-zero model through $q_n$ fails.
The integer $d$ is defined exactly below; its numerical value and its
possible odd prime divisors are not computed here.
\end{corollary}
\begin{proof}
Lift each binary certificate coefficient to its integer representative
$0$ or $1$, and let $\widetilde d_i$ denote the integer Fox images. Put
$S=\widetilde X_0\widetilde d_0+\widetilde X_1\widetilde d_1$.
Equation~\eqref{eq:a8unit} gives $S=1+2W$ for some $W\in\Z[A_8]$.
On the free abelian group $\Z[A_8]$ of rank $20160$, consider right
multiplication $T_S:a\mapsto aS$. Its integer matrix reduces to the
identity modulo two. Thus $d=\det(T_S)$ is odd and nonzero.
Apply $T_S\operatorname{adj}(T_S)=d\,\mathrm{id}$ to the identity
basis element. With $c=\operatorname{adj}(T_S)(1)$ this gives
$cS=d\cdot1$, so the image ideal contains $d\cdot1$.
A unital model into $M_m(R)$ killing the ideal forces
$dI_m=0$, hence $d1_R=0$. A nonzero ring then has positive
characteristic dividing $|d|$. The argument uses the exact oddness
of the determinant, not an unperformed numerical determinant
calculation. It does not establish $d=\pm1$.
\end{proof}

\begin{corollary}\label{cor:lifts}
The exclusion in Theorem~\ref{thm:a8} holds for every model through
$q_n$ over a nonzero commutative coefficient ring in which $2$ is
nilpotent, including every ring of characteristic $2^r$.
More generally, there is no unital matrix model killing $I_n$ over
a commutative local ring whose residue field has characteristic two
and whose residue group action factors through $q_n$.
\end{corollary}
\begin{proof}
Lift the two certificate coefficients to integer group-ring elements.
The lifted identity has the form $1+2W$ in $\Z[A_8]$ and belongs to
the ideal generated by the lifted first-row derivatives.
If $2^r=0$, then $(2W)^r=0$; the finite geometric series makes
$1+2W$ invertible. Hence the image ideal over the coefficient ring
contains one. For the local statement, reduce the proposed matrix
model modulo its maximal ideal. It remains unital on a nonzero
matrix ring over the residue field, factors through $q_n$ by
hypothesis, and kills all derivatives, contradicting
Theorem~\ref{thm:a8}. The action before reduction need not itself
factor through $q_n$.
\end{proof}

\section{Exact replay and remaining rank problem}
From this directory, run
\begin{quote}\ttfamily
python3 verify\_uniform\_fox\_ideal\_models.py
\end{quote}
The frozen companion input has SHA-256
\url{29464dbf71101d8373eecc89da9065261e4689580bd741313afa240e4c7fa25b};
the permutation-tuple certificate has SHA-256
\url{a070353a85519beb796cc24773f6ea06fadd760005352914a0f1fc0760a7f3da}.
Both files are retained beside this manuscript.

The verifier uses Python's standard library. It checks the hashes of
the frozen knot and signed-diagram inputs, reconstructs the literal
free words, binds all seventeen crossings and verifies the redundant
closure relation through ordered-product preservation. It computes
the scalar first row and checks \eqref{eq:scalarunit} by integer
convolution. It independently propagates the permutations
\eqref{eq:top}, checks the six-step return and enumerates their group.
It computes the four first-row derivatives as integer permutation
sums, checks the Fox fundamental identity, and verifies
\eqref{eq:a8unit} by exact sparse convolution modulo two.
Corrupted coefficient, index and input controls are rejected.
Finite tail evaluations are controls for the written induction, rather
than its proof. The search that discovered the coefficients is outside
the logical dependency chain of their verification.

The whole noncommutative ideal $I_n\subset\Z[G_n]$ has not been
shown to contain one. Other quotient kernels, the possible odd exceptional characteristics
for this fixed quotient, characteristic-zero noncommuting models
through other quotients, and other generating systems remain
unclassified. These results imply neither
$\operatorname{rank}(G_n)\le3$ nor $\operatorname{rank}(G_n)=4$.
The already proved meridional rank four remains independent.
The Moriah theorem is sufficient; failure of these models is not a
converse rank test. Meridional tameness and worldwide priority are
separate open comparisons. All results here retain internal exact
replay status, with formal verification and external specialist review
unset.

\paragraph{AI assistance and review.}
ChatGPT (OpenAI) assisted with mathematical arguments, literature checks,
exposition, proof review and verification code.

\begin{thebibliography}{9}
\bibitem{moriah} Y.~Moriah,
\emph{Heegaard splittings of knot exteriors},
Geometry \& Topology Monographs \textbf{12} (2007), 191--232.
\href{https://doi.org/10.2140/gtm.2007.12.191}{doi:10.2140/gtm.2007.12.191}.
Adopted scope: setup on p.~197 and Theorem~3.5 on p.~198;
original upstream lemma remains imported.
\bibitem{family} Accompanying exact-family manuscript,
\emph{An infinite prime four-bridge family beyond augmentation and
Coxeter bounds}, \texttt{PRIME\_AUGMENTATION\_GAP\_FAMILY.tex},
Theorem \texttt{thm:main}. Used only for the established motivating
family status, rather than either new unit identity.
\end{thebibliography}
\end{document}
